\chapter{Partikeln}
\label{cha:Partikeln}

Partikeln sind kleine Wörter, die die Stimmung, Einstellung oder Nuancen des Sprechers ausdrücken. Sie sind besonders wichtig für die mündliche Kommunikation.

\section{Abtönungspartikeln (Modalpartikeln)}
\label{sec:Modalpartikeln}

Abtönungspartikeln verändern die Bedeutung eines Satzes subtil:

\begin{center}
\begin{tabular}{|l|c|l|}
\hline
\textbf{Partikel} & \textbf{Funktion} & \textbf{Beispiel} \\ \hline
aber & Erstaunen, Bewertung & Das ist aber schön! \\ \hline
also & Schlussfolgerung & Also, das war's. \\ \hline
auch & Einbeziehung & Das kann ich auch. \\ \hline
bloß & Einschränkung & Beeil dich bloß! \\ \hline
denn & Frage & Wie geht's denn? \\ \hline
doch & Widerspruch/Erstaunen & Das ist doch klar! \\ \hline
eben & Resignation & Das ist eben so. \\ \hline
eigentlich & Einschränkung & Was meinst du eigentlich? \\ \hline
etwa & Frage & Ist das etwa wahr? \\ \hline
ja & Zustimmung & Das ist ja klar. \\ \hline
mal & Abschwächung & Komm mal vorbei! \\ \hline
nur & Einschränkung & Nur zu! \\ \hline
schon & Beruhigung & Das ist schon okay. \\ \hline
toll & Begeisterung & Das ist ja toll! \\ \hline
vielleicht & Erstaunen & Das ist vielleicht schwer! \\ \hline
wohl & Vermutung & Das wird wohl stimmen. \\ \hline
\end{tabular}
\end{center}

\section{Gradpartikeln}
\label{sec:Gradpartikeln}

Gradpartikeln verstärken oder relativieren die Bedeutung:

\begin{center}
\begin{tabular}{|l|c|l|}
\hline
\textbf{Partikel} & \textbf{Funktion} & \textbf{Beispiel} \\ \hline
sehr & Verstärkung & Das ist sehr gut. \\ \hline
ziemlich & Annäherung & Das ist ziemlich gut. \\ \hline
recht & Annäherung & Das ist recht gut. \\ \hline
ganz & Annäherung & Das ist ganz gut. \\ \hline
ein bisschen & Abschwächung & Ein bisschen besser. \\ \hline
ein wenig & Abschwächung & Ein wenig besser. \\ \hline
bisschen & Abschwächung & Das ist bisschen schwer. \\ \hline
\end{tabular}
\end{center}

\section{Steigerungspartikeln}
\label{sec:Steigerungspartikeln}

\begin{center}
\begin{tabular}{|l|c|l|}
\hline
\textbf{Partikel} & \textbf{Funktion} & \textbf{Beispiel} \\ \hline
noch & Steigerung & Das ist noch besser. \\ \hline
schon & Vorzeitigkeit & Das war schon gut. \\ \hline
immer & Wiederholung & Immer besser. \\ \hline
\end{tabular}
\end{center}

\section{negationspartikeln}
\label{sec:Negationspartikeln}

\begin{center}
\begin{tabular}{|l|c|l|}
\hline
\textbf{Partikel} & \textbf{Funktion} & \textbf{Beispiel} \\ \hline
nicht & Negation & Das ist nicht gut. \\ \hline
kein & Negation & Das ist kein Problem. \\ \hline
nie & Negation & Das war nie einfach. \\ \hline
niemals & starke Negation & Das ist niemals einfach. \\ \hline
keineswegs & starke Negation & Das ist keineswegs einfach. \\ \hline
\end{tabular}
\end{center}

\section{Antwortpartikeln}
\label{sec:Antwortpartikeln}

\begin{center}
\begin{tabular}{|l|c|l|}
\hline
\textbf{Partikel} & \textbf{Funktion} & \textbf{Beispiel} \\ \hline
ja & Zustimmung & Ja, das stimmt. \\ \hline
nein & Ablehnung & Nein, das stimmt nicht. \\ \hline
doch & Widerspruch & Nein, das ist doch falsch! \\ \hline
doch & positive Antwort & Hast du Zeit? - Doch. \\ \hline
vielleicht & Unsicherheit & Vielleicht. \\ \hline
na ja & skeptisch & Na ja, maybe. \\ \hline
\end{tabular}
\end{center}

\section{Übungen zu Partikeln}
\label{sec:PartikelnUebungen}

\begin{enumerate}
    \item Komm \underline{\hspace{1cm}} vorbei! (mal)
    \item Das ist \underline{\hspace{1cm}} schön! (aber)
    \item Wie geht's \underline{\hspace{1cm}}? (denn)
    \item Das ist \underline{\hspace{1cm}} klar! (doch)
    \item Beeil dich \underline{\hspace{1cm}}! (bloß)
\end{enumerate}

\chapterend